\documentclass[10pt]{article}

\usepackage[utf8]{inputenc}

\usepackage[T1]{fontenc}

\usepackage{amsmath}

\usepackage{amsfonts}

\usepackage{amssymb}

\usepackage[version=4]{mhchem}

\usepackage{stmaryrd}

\usepackage{bbold}

\usepackage{graphicx}

\usepackage[export]{adjustbox}

\graphicspath{ {./images/} }



\begin{document}

КООРДИНАЦИОННОЕ УПОРЯДОЧЕНИЕ - сМ. ДальНиЙ и ближний порядок.



КОПРИСОЕДИНЕННОЕ ПРЕДСТАВЛЕНИЕ - представление группы Ли $G$, контрагредиентное к присоединенному представлению ad этой группы.



Пусть $G$ - группа Ли, $\mathbb{S}=T_{e} G$ - ее алгебра Ли. Группа Ли $G$ действует на себя сопряжением. Дифференциал этого действия в единице $e \in G$ определяет присоединенное дей-



\includegraphics[max width=\textwidth, center]{2023_07_08_44290df9b3a193afefb0g-1(1)}

женное действие $\mathrm{ad}^{*}: G \rightarrow G L\left(\right.$ (ङ) $\left.^{*}\right)$ и называется к о п р и с о единенным п ред ставлен ием. Соответствующее K. п. алгебры Ли ad*: (\$) $\rightarrow$ End ((ङ) $)$ в явном виде задается формулой



\[

\operatorname{ad}_{x}^{*} \xi=\xi([\cdot, x]), x \in \mathbb{S}, \xi \in \mathbb{S}^{*} .

\]



Коалгебра Ли (s ${ }^{*}$ имеет естественную пуассонову структуру, инвариантную относительно действия группы $G$ : скобка Ли на элементах \& - линейных функциях на (\$) по правилу Лейбница продолжается до скобки Пуассона на всех функциях на (s**. Симплектич. слои этой пуассоновой структуры суть орбиты коприсоединенного действия. Соответствующая симплектич. структура на орбитах задается Ли-Березина-Костанта-Кириллова формой. М. Э. Казарян.



КОРАБЕЛЬНЫЕ УГЛЫ - см. Ориентации уалы.



KOPАЗМЕРНОСТЬ векторного пространствасм. Векторное пространство.



KOPAЗMEPHОСТЬ $m$-мерного подмногообраз и я $Y^{m}$ в $n$-мерном многообразии $X^{n}-$ разность $n-m$.



КОРАЗМЕРНОСТЬ с ло е н и я - см. Слоение.



КОРИОЛИСА СИЛА - одна из сил инерции; вводится для учета влияния вращения подвижной системы отсчета на относительное движение материальной точки; названа по имени Г. Кориолиса (G. Coriolis). K. c. равна произведению массы точки на ее Кориолиса ускорение и направлена противоположно этому ускорению. Эффект, учитываемый К. с., состоит в том, что во вращающейся системе отсчета материальная точка, движущаяся непараллельно оси этого вращения, отклоняется по направлению, перпендикулярному к ее относительной скорости, или оказывает давление на тело, препятствующее такому отклонению. На Земле этот эффект, обусловленный ее суточным вращением, заключается в том, что свободно падающие тела отклоняются от вертикали к востоку (в первом приближении), а тела, движущиеся вдоль земной поверхности, отклоняются в Северном полушарии вправо, а в Южном - влево от направления их движения. Вследствие медленного вращения Земли эти отклонения весьма малы и заметно сказываются или при очень больших скоростях движения (напр., у ракет и у артиллерийских снарядов с большими дальностями полета), или когда движение длится очень долго, напр. подмыв соответствующих берегов рек (так наз. закон Бэра), возникновение нек-рых воздушных и морских течений и др. В технике К. с. учитывается в теории гироскопов, турбин и многих других вращающихся систем.



C.M. Tape.



KОРИОЛИСА УСКОРЕНИЕ, поворотно е ускоре н и е, - составляющая полного ускорения точки, которая появляется при так наз. сло жном д в и жени и, когда переносное движение, то есть движение подвижной системы отсчета, не является поступательным. К. у. возникает вследствие изменения относительной скорости точки $v_{\text {от }}$ при переносном движении и переносной скорости при относительном движении точки. Численно К. у. равно:



\[

w_{\text {кор }}=2 \omega_{\text {пер }} v_{\text {от }} \sin \alpha ;

\]



как вектор К. у. определяется формулой



\[

\boldsymbol{w}_{\text {кор }}=2\left[\omega_{\text {пер }} v_{\text {от }}\right] \text {, }

\]



где $\omega_{\text {пер }}-$ угловая скорость поворота подвижной системы отсчета относительно неподвижной, $\alpha-$ угол между $v_{\text {от }}$ и $\omega_{\text {пер. }}$ Направление К.у. можно получить, спроектировав



\includegraphics[max width=\textwidth, center]{2023_07_08_44290df9b3a193afefb0g-1}

вернув эту проекцию на $90^{\circ}$ в сторону переносного вращения. Таким образом, К. у.- это часть ускорения точки по отношению к основной, а не к подвижной системе отсчета. Напр., при движении вдоль поверхности Земли вследствие ее вращения точка будет иметь К. у. по отношению к звездам, а не к Земле. К.у. равно нулю при поступательном переносном движении $\left(\omega_{\text {пер }}=0\right)$ или когда $\alpha=0$. С. М. Тарг. KОРНЕВОЕ ПОДПРОСТРАНСТВО - см. Вec.



КОРОТКИЕ ВОЛНЫ в анизотропных средахволны с длиной волны $\lambda \ll L$, где $L-$ масштаб неоднородности анизотропной среды. В физике наиболее часто встречаются упругие, диэлектрические и магнитоактивные плазменные анизотропные среды. Если $c_{i k j l}(M)-$ тензор упругих постоянных, $\varepsilon_{i k}(M)$ - тензор диэлектрич. проницаемости, $M\left(x_{1}, x_{2}, x_{3}\right)$ - точка пространства, то обычно полагают



\[

L=\min _{i, k, j, l}\left(c_{i k j l}\left|\operatorname{grad} c_{i k j l}\right|^{-1}\right) ; \quad L=\min _{i, k}\left(\varepsilon_{i k}\left|\operatorname{grad} \varepsilon_{i k}\right|^{-1}\right) .

\]



К. в. в анизотропных средах локально в окрестности каждой фиксированной точки $M$ представляют собою плоские волны и могут быть представлены разложениями по обратным степеням частоты $\omega$ :



\[

\begin{aligned}

\boldsymbol{U} & =\boldsymbol{A}(\boldsymbol{M}, \omega) \exp [i \omega(\tau(\boldsymbol{M})-\boldsymbol{t})], \\

A(M, \omega) & =\sum_{n=0}^{\infty} \boldsymbol{A}^{(n)}(M)(-i \omega)^{-n}, \quad A^{(0)} \neq 0 .

\end{aligned}

\]



Здесь $U-$ вектор смешений частиц анизотропной среды или шестимерный вектор $\boldsymbol{U}=(\boldsymbol{E}, \boldsymbol{H})$ напряженности электромагнитного поля, $t$ - время, $\tau(M)$ - эйконал и $A^{(n)}(M)-$ коэффициент разложения амплитуды $\boldsymbol{A}(M, \omega)$. В должным образом введенных безразмерных переменных разложение (1) - разложение по малому безразмерному параметру $x=(\lambda / L)_{0}$, где $(\lambda / L)_{0}-$ значение отношения $\lambda / L$ в нек-рой фиксированной точке $M_{0}$.



Подстановка $U$ в однородные уравнения движения анизотропной среды и приравнивание нулю коэффициентов при различных степенях $1 /(-i \omega)$ дает для составляющих $\boldsymbol{A}^{(0)}$ однородную и для $\boldsymbol{A}^{(n)}, n \geqslant 1$, неоднородную рекуррентную систему уравнений. Однородная система для $\boldsymbol{A}^{(0)}$ имеет отличные от нуля решения, если ее определитель равен нулю. Напр., в случае упругой анизотропной среды $\operatorname{det}\left(C_{i j}-\delta_{i j}\right)=0$, где $C_{i j}=\rho^{-1} c_{i k j l} \tau_{l} \tau_{k}, \tau_{l}=\partial \tau / \partial x_{l}, \rho-$ плотность. Матрица $C_{i j}$ является положительно определенной, еe собственные числа $\lambda_{r}\left(\tau_{l}\right)>0, r=1,2,3$,- однородные второй степени функции $\tau_{l}-$ при любых $\tau_{l}$ удовлетворяют неравенствам $\lambda_{1}>\lambda_{2} \geqslant \lambda_{3}$. Определитель будет равен нулю, только если $\lambda_{r}\left(\tau_{l}\right)=1$ при каком-либо $r$. Уравнение $\lambda_{r}\left(\tau_{l}\right)=1-$ уравнение эйконала для функции $\tau(M)$. В неоднородной анизотропной среде могут распространяться три типа коротких волн: квазипродольная волна $(r=1)$ и две квазипоперечные $(r=2,3)$. Фронт волны типа $r$ распространяется с фазовой скоростью



\[

\boldsymbol{v}_{r}^{\Phi}=\omega \sqrt{\lambda_{r}\left(n_{l}\right)} \boldsymbol{n}

\]



в направлении нормали к фронту



\[

\boldsymbol{n}=\operatorname{grad} \tau /|\operatorname{grad} \tau|

\]



Лучам и в анизотропных средах называются векторные линии поля векторов групповой скорости



\[

\boldsymbol{v}_{r}^{\mathrm{rp}}=\partial v_{r}^{\Phi} / \partial \boldsymbol{n} .

\]



Энергия волны в нулевом (геометро-оптическом) приближении распространяется по лучевым трубкам с групповой скоростью. В анизотропных средах в отличие от изотропных лучи неортогональны волновым фронтам.





\end{document}